\chapter{A Concrete Naming Certificate for $\BishopUniformContinuityModulusUp$}
\label{ch:concrete-instances-bishop-uniform-continuity-modulus-namecert}
\origin{human}

Following Bishop and Bridges, Constructive Analysis, the $\BishopUniformContinuityModulusUp$ packet records an explicit tolerance-to-source-window modulus as finite BEDC NameCert data. It sits on the finite stream surface of \autoref{ch:concrete-instances-streamname-namecert}, the regular rational readback of \autoref{ch:concrete-instances-regseqrat-namecert}, the terminal real-name seal of \autoref{ch:concrete-instances-real-namecert}, the metric source rows of \autoref{ch:concrete-instances-metric-namecert}, the continuous-map graph of \autoref{ch:concrete-instances-continuousmap-namecert}, and the downstream modulus consumer of \autoref{ch:concrete-instances-uniformmodulus-namecert}. The packet names the constructive uniform-continuity modulus needed by the open L10 real and completion route; it does not import open-cover compactness, countable choice, quotient real equality, selected representatives, or ambient $\RealUp$ completeness.

\begin{definition}[Bishop uniform-continuity modulus carrier]
\label{def:bishop-uniform-continuity-modulus-carrier}
A $\BishopUniformContinuityModulusUp$ carrier is a finite $\BHist$ packet
$$
M=(S,F,T,L,D,R,E,H,C,P,N)
$$
with the following displayed rows:
\begin{enumerate}
\item $S$ is the $\StreamNameUp$ source-window row naming the finite input observations that a tolerance request may inspect.
\item $F$ is the $\ContinuousMapUp$ graph row, including the displayed pointwise continuity data whose local bounds are eligible for uniform readback.
\item $T$ is the tolerance-row carrier admission ledger, recording the requested target precision and the finite source-window budget produced for that request.
\item $L$ is the pointwise-to-uniform modulus ledger, displaying the finite route from local continuity rows to one uniform source tolerance.
\item $D$ is the dyadic lower-bound readback row, computed inside the $\DyadicRatCoreUp$ tolerance discipline and exposed before any real-valued consumer reads the modulus.
\item $R$ is the $\RegSeqRatUp$ readback row attached to the source and target estimates.
\item $E$ is the $\RealUp$ seal row, reached only after $S,F,T,L,D,R$ have been displayed.
\item $H$ records componentwise $\hsame$ transport for source-window, continuous-map, tolerance, ledger, dyadic, regular-readback, real-seal, replay, provenance, and naming rows.
\item $C$ records displayed $\Cont$ replay from a target-tolerance request through the same finite modulus route.
\item $P$ is the $\Pkg$ provenance over $\BishopUniformContinuityModulusUp$, $\StreamNameUp$, $\ContinuousMapUp$, $\UniformModulusUp$, $\DyadicRatCoreUp$, $\RegSeqRatUp$, $\RealUp$, $\BHist$, $\hsame$, $\Cont$, and $\NameCert$ rows.
\item $N$ is the local $\NameCert_{\BishopUniformContinuityModulusUp}$ row.
\end{enumerate}
The classifier compares accepted packets only through $S,F,T,L,D,R,E,H,C,P,N$. It has no coordinate for open-cover compactness, countable choice, selected representatives, quotient equality of real names, a hidden host modulus, or ambient completion.
\end{definition}

\begin{theorem}[Bishop uniform-continuity modulus NameCert obligations]
\label{thm:bishop-uniform-continuity-modulus-namecert-obligations}
For every accepted $\BishopUniformContinuityModulusUp$ carrier $M=(S,F,T,L,D,R,E,H,C,P,N)$, the local $\NameCert_{\BishopUniformContinuityModulusUp}$ surface consists exactly of the following five obligation rows:
\begin{enumerate}
\item tolerance-row carrier admission through $T$;
\item pointwise-to-uniform modulus ledger exposure through $L$;
\item dyadic lower-bound readback through $D$;
\item $\RealUp$ seal nonescape through $E$;
\item provenance and name exhaustion through $P,N$ together with transport and replay rows $H,C$.
\end{enumerate}
Every $\StreamNameUp$, $\RegSeqRatUp$, $\UniformModulusUp$, or $\RealUp$ consumer of the packet must read those rows through the displayed carrier route, and no accepted read may project compactness, choice, quotient equality, selected representatives, or a host modulus.
\end{theorem}
\begin{proof}
Project the carrier of \autoref{def:bishop-uniform-continuity-modulus-carrier}. The source-window row $S$ and continuous-map row $F$ are visible before any tolerance output is read.

The tolerance row $T$ and ledger row $L$ record the constructive step from a requested target precision to a uniform source window. The dyadic row $D$ supplies the finite lower-bound comparison, and the regular readback row $R$ makes that comparison available to real-name consumers.

Finally the seal row $E$ is downstream of the displayed route, while $H,C,P,N$ transport, replay, source, and name the same finite packet. Since no other carrier coordinate exists, the listed five obligation rows exhaust the local NameCert surface.
\end{proof}

\begin{theorem}[Bishop uniform-continuity modulus compact handoff]
\label{thm:bishop-uniform-continuity-modulus-compact-handoff}
Any compact-uniform-continuity, continuous-map, or $\RealUp$ consumer that reads an accepted $\BishopUniformContinuityModulusUp$ packet must factor its read through
$$
\StreamNameUp,\ \ContinuousMapUp,\ \DyadicRatCoreUp,\ \RegSeqRatUp,\ \UniformModulusUp,\ \RealUp,\ \BHist,\ \hsame,\ \Cont,\ \Pkg,\ \NameCert .
$$
Thus the packet supplies a finite constructive modulus object for later compact-uniform-continuity and completion-facing chapters without importing choice, open-cover compactness, selected representatives, quotient equality, or ambient real completeness.
\end{theorem}
\begin{proof}
Use \autoref{thm:bishop-uniform-continuity-modulus-namecert-obligations}. A downstream read first exposes the finite source window, the continuous-map row, the tolerance row, the pointwise-to-uniform ledger, and the dyadic lower-bound readback.

The regular rational row and real seal then pass the same displayed estimates to $\RegSeqRatUp$ and $\RealUp$ consumers. The $\UniformModulusUp$ output is therefore a named row of the finite packet, not a hidden host function.

Transport, replay, provenance, and local naming preserve only the displayed coordinates. The compact handoff can cite the modulus object, but it receives no open-cover compactness, choice principle, quotient equality, selected representative, or ambient completeness theorem from this chapter.
\end{proof}

\closureat{BishopUniformContinuityModulusUp}{seedStr}

\begin{closurestatus}{\BishopUniformContinuityModulusUp}
  \constructivestory{A $\BishopUniformContinuityModulusUp$ packet is generated from source-window rows, continuous-map rows, tolerance rows, pointwise-to-uniform ledgers, dyadic lower-bound readback, regular rational readback, a Real seal, componentwise hsame transport, displayed Cont replay, Pkg provenance, and a local NameCert row.}
  \theoryclosure{\seedClosure}
  \scopeclosed{\autoref{def:bishop-uniform-continuity-modulus-carrier}, \autoref{thm:bishop-uniform-continuity-modulus-namecert-obligations}, and \autoref{thm:bishop-uniform-continuity-modulus-compact-handoff} seed the finite Bishop uniform-continuity modulus route for L10 real and completion-facing consumers.}
  \formalstatus{\unformalizedV}
  \bridgestatus{none}
  \notclaimed{This chapter does not close open-cover compactness, countable choice, selected representatives, quotient real equality, host modulus functions, ambient $\RealUp$ completeness, Lean verification, or bridge checking.}
  \upgradepath{Obligation closure requires checked carrier admission, tolerance-row admission, pointwise-to-uniform ledger exposure, dyadic lower-bound readback, Real seal nonescape, and provenance/name exhaustion for BEDC.Derived.BishopUniformContinuityModulusUp.}
\end{closurestatus}
